\chapter{Base de Scrap e Análise de Ocorrências}

\label{cap:manual-4}

\section{Consulta, Ordenação e Filtros da Base de Scrap}

\subsection*{Objetivo}

Consultar todas as ocorrências de refugo importadas para o sistema, permitindo ordenar por data ou valor, filtrar por status de análise e localizar rapidamente os itens que precisam de investigação técnica.\par

\subsection*{Antes de começar}

\begin{itemize}
\item Estar conectado ao sistema e acessar o menu lateral \textbf{Base de Scrap}.
\end{itemize}

\begin{remark}[info]
A abertura da Base de Scrap e seus controles de ordenação e status já estão descritos no guia rápido. Consulte o passo~\ref{passo:5} (\pagina{passo:5}).
\end{remark}
\subsection*{Opções de consulta}
\begin{itemize}
\item A ordenação aceita \textbf{Data da transação}, \textbf{Organização}, \textbf{Código do item}, \textbf{Quantidade} e \textbf{Valor (USD)}, em ordem crescente ou decrescente.
\item O status \textbf{Não analisado} identifica ocorrências aguardando análise; \textbf{Rascunho}, análises iniciadas e salvas parcialmente; \textbf{Revisado}, análises concluídas.
\end{itemize}


\passo[0.12000/0.37547/0.04706/0.28302,0.44941/0.25849/0.37647/0.16981,0.79529/0.25849/0.87059/0.16981,0.44941/0.43774/0.37647/0.53774,0.79529/0.43774/0.87059/0.53774,0.44941/0.61887/0.37647/0.70755,0.92000/0.92453/0.91765/0.83019,0.79765/0.92453/0.75294/0.83019]{completo/16-base-scrap-filtros}{Aplicar filtros da base}{%
Abra Filtros da base pelo botão Filtrar. O detalhe mostra a categoria Produção.\par
\begin{elementos}
\item \textbf{Categorias.} Alterne Geral, Catálogo, Produção e Análise.
\item \textbf{Produção.} Selecione o produto ou o grupo de produção.
\item \textbf{Departamento.} Restrinja o recorte ao departamento desejado.
\item \textbf{Divisão.} Escolha a divisão, quando necessário.
\item \textbf{Filtro da linha.} Defina o agrupamento de linhas.
\item \textbf{Linha.} Escolha a linha de produção.
\item \textbf{Aplicar.} Atualize a tabela com o recorte selecionado.
\item \textbf{Limpar filtros.} Remova as restrições e volte à lista completa.
\end{elementos}
}

\Needspace{160pt}

\subsection*{Campos e regras}

\begingroup
\fontsize{9}{12}\selectfont
\setlength{\LTpre}{4pt}
\setlength{\LTpost}{8pt}
\begin{longtable}{@{}L{3.389cm}L{1.916cm}L{2.652cm}L{6.778cm}@{}}
\toprule
\guiacabtabela \thead{Campo} & \thead{Obrigatório} & \thead{Formato} & \thead{Descrição e regra de preenchimento} \\
\midrule
\endfirsthead
\toprule
\guiacabtabela \thead{Campo} & \thead{Obrigatório} & \thead{Formato} & \thead{Descrição e regra de preenchimento} \\
\midrule
\endhead
\bottomrule
\endlastfoot
\rowcolor{white} \textbf{Ordenar por} & Não & Campo de seleção & Organiza a tabela por \tech{Data da transação}, \tech{Organização}, \tech{Código do item}, \tech{Quantidade} ou \tech{Valor (USD)}. \\
\rowcolor{guiaareia!70} \textbf{Status da análise} & Não & Campo de seleção & Filtra entre \tech{Todos}, \tech{Não analisado}, \tech{Rascunho} e \tech{Revisado}. \\
\rowcolor{white} \textbf{Ordem (``Decrescente'' / ``Crescente'')} & Não & Botão alternador & Inverte a listagem do maior para o menor ou do mais recente para o mais antigo. \\
\rowcolor{guiaareia!70} \textbf{Filtros da base (``Geral'', ``Catálogo'', ``Produção'', ``Análise'')} & Não & Painel de filtros & Permite restringir a consulta por linha de produção, departamento, item ou analista responsável. \\
\end{longtable}
\endgroup

\subsection*{Mensagens e solução de problemas}

\begin{remark}
\textbf{Algumas ocorrências não aparecem na tabela}\par
\textbf{Causa provável:} Há um filtro ativo no painel \textbf{Filtrar} ou o seletor de status não está marcado em \textbf{Todos}.\par
\textbf{Como resolver:} Clique em \textbf{Filtrar > Limpar filtros} e selecione \textbf{Todos} no filtro de status da análise.
\end{remark}

\section{Preenchimento e Finalização de Análise de Scrap Individual}

\subsection*{Objetivo}

Registrar a investigação técnica de uma ocorrência de scrap em 7 etapas (O problema, Onde e quando, Causa, Risco e responsabilidade, Custo, Contenção imediata e Evidências/fotos), permitindo salvar como rascunho ou finalizar a análise.\par

\subsection*{Antes de começar}

\begin{itemize}
\item Categorias de scrap já cadastradas em \textbf{Configurações > Tipos de Scrap}.
\end{itemize}

\passo[0.46154/0.26554/0.21154/0.23164,0.46154/0.40113/0.21154/0.46893,0.15385/0.65537/0.09615/0.61582,0.38462/0.65537/0.40385/0.60452,0.61538/0.65537/0.67308/0.60452,0.24038/0.75367/0.11538/0.77401,0.65385/0.75367/0.75000/0.77401]{completo/17-ocorrencia-problema-local}{Descrever o problema e o local}{%
Na Base de Scrap, abra uma ocorrência Não analisada ou em Rascunho para preencher a análise.\par
\begin{elementos}
\item \textbf{Título da análise.} Resuma o defeito ou a causa.
\item \textbf{Descrição e causa raiz.} Explique o ocorrido, as causas identificadas e as ações tomadas. O campo aceita até 20.000 caracteres.
\item \textbf{Data.} Informe a data da ocorrência.
\item \textbf{Hora.} Informe o horário.
\item \textbf{Turno.} Selecione o turno.
\item \textbf{Linha.} Informe a linha.
\item \textbf{Configuração da linha.} Complete a configuração correspondente.
\end{elementos}
}

\passo[0.30808/0.19459/0.11616/0.15135,0.50000/0.29730/0.50000/0.23784,0.29798/0.41622/0.11616/0.45405,0.82323/0.37297/0.89394/0.32432,0.50000/0.62162/0.50000/0.56216,0.24747/0.74595/0.14646/0.70270,0.81313/0.74595/0.89394/0.70270,0.87374/0.96973/0.80303/0.90811,0.96667/0.96973/0.95455/0.90811]{completo/18-ocorrencia-causa-risco}{Registrar causa e responsabilidade}{%
\begin{elementos}
\item \textbf{Tipo de scrap.} Escolha a categoria da ocorrência.
\item \textbf{Categoria 4M.} Selecione Mão de obra, Máquina, Material, Método ou Outro.
\item \textbf{Causa analisada.} Descreva a causa investigada.
\item \textbf{Situação da causa.} Informe se é hipótese ou causa confirmada.
\item \textbf{Nível de risco.} Escolha Não avaliado, Baixo, Moderado, Alto ou Crítico.
\item \textbf{Onde/como.} Registre a segmentação do problema, por exemplo Empilhamento na saída da linha.
\item \textbf{Área responsável.} Informe a área responsável pela ocorrência.
\item \textbf{Salvar rascunho.} Guarde o preenchimento e continue depois.
\item \textbf{Finalizar relatório.} Conclua a análise quando todos os dados e evidências estiverem prontos.
\end{elementos}
\par\smallskip {\fontsize{8.2}{11}\selectfont\color{guiacinza}Antes de finalizar, role até Custo, Contenção imediata e Evidências e fotos. Confira Tipo de custo, Modelo e Quantidade. Se Quantidade ficar em branco, será usada a quantidade original da base. Registre a contenção e anexe até oito imagens JPEG, PNG ou WebP, com até 10 MB por foto.\par}
}

\passo[0.76863/0.75102/0.86275/0.73469]{completo/19-finalizar-analise}{Finalizar a análise}{
\begin{elementos}
\item \textbf{Finalizar relatório.} Na janela Finalizar análise, confira o aviso e clique em Finalizar relatório para confirmar. Após a finalização, a análise continuará editável apenas por seu autor.
\end{elementos}
}

\begin{remark}[info]
A consulta de uma ocorrência revisada já está descrita no guia rápido. Consulte o passo~\ref{passo:6} (\pagina{passo:6}).
\end{remark}
Depois de finalizar a análise, a etiqueta \textbf{Revisado} confirma a conclusão.
O ícone de lápis permite editar a análise; o coração abre o cadastro de modelo favorito,
detalhado na seção seguinte.


\Needspace{160pt}

\subsection*{Campos e regras}

\begingroup
\fontsize{9}{12}\selectfont
\setlength{\LTpre}{4pt}
\setlength{\LTpost}{8pt}
\begin{longtable}{@{}L{3.389cm}L{1.916cm}L{2.652cm}L{6.778cm}@{}}
\toprule
\guiacabtabela \thead{Campo} & \thead{Obrigatório} & \thead{Formato} & \thead{Descrição e regra de preenchimento} \\
\midrule
\endfirsthead
\toprule
\guiacabtabela \thead{Campo} & \thead{Obrigatório} & \thead{Formato} & \thead{Descrição e regra de preenchimento} \\
\midrule
\endhead
\bottomrule
\endlastfoot
\rowcolor{white} \textbf{Título da análise *} & Sim & Texto curto & Resumo breve que identifica a análise nas listas e relatórios. \\
\rowcolor{guiaareia!70} \textbf{Descrição e causa raiz *} & Sim & Texto longo (até 20.000 caracteres) & Detalhamento completo da ocorrência, causas e ações imediatas. \\
\rowcolor{white} \textbf{Data / Hora / Turno / Linha} & Não & Data (\tech{AAAA/MM/DD}), Hora e Seletores & Identificam o momento e o local da linha onde ocorreu a perda. \\
\rowcolor{guiaareia!70} \textbf{Tipo de scrap *} & Sim & Campo de seleção & Categoria escolhida dentre os \textbf{Tipos de Scrap} ativos. \\
\rowcolor{white} \textbf{Categoria 4M da causa} & Não & Botões de opção (\tech{Mão de obra}, \tech{Máquina}, \tech{Material}, \tech{Método}, \tech{Outro}) & Classifica a natureza da causa para análise nos relatórios. \\
\rowcolor{guiaareia!70} \textbf{Situação da causa / Nível de risco} & Não & Seletor e Botões de opção & Indicam se a causa ainda é uma hipótese e qual a gravidade do risco (\tech{Baixo} a \tech{Crítico}). \\
\rowcolor{white} \textbf{Onde/como (segmentação) / Área responsável} & Não & Texto livre & Detalham o posto ou etapa da linha e o setor responsável pela tratativa. \\
\rowcolor{guiaareia!70} \textbf{Tipo de custo (IF Cost) / Quantidade} & Não & Seletor e Número & Classificam o custo; quando a quantidade fica vazia, o sistema adota a quantidade da base. \\
\rowcolor{white} \textbf{Contenção imediata} & Não & Texto livre & Ação imediata executada no turno para interromper a perda. \\
\rowcolor{guiaareia!70} \textbf{Evidências e fotos (``0/8'')} & Não & Envio de imagens (\tech{JPEG}, \tech{PNG}, \tech{WebP}) & Permite anexar até \textbf{8 fotos} de no máximo \textbf{10 MB} cada. \\
\end{longtable}
\endgroup

\subsection*{Mensagens e solução de problemas}

\begin{remark}
\textbf{\textbullet{} Alterações não salvas (no canto inferior esquerdo)}\par
\textbf{Causa provável:} Você alterou informações na tela, mas ainda não clicou em \textbf{Salvar rascunho} nem em \textbf{Finalizar relatório}.\par
\textbf{Como resolver:} Clique em \textbf{Salvar rascunho} ou \textbf{Finalizar relatório} antes de fechar a janela.
\end{remark}

\begin{remark}
\textbf{O relatório será finalizado e continuará editável apenas por quem o criou. Deseja continuar?}\par
\textbf{Causa provável:} Aviso de segurança informando que, após finalizar, apenas o autor da análise poderá editá-la.\par
\textbf{Como resolver:} Clique em \textbf{Finalizar relatório} para confirmar. Caso outro colega ainda precise complementar dados antes, cancele e use \textbf{Salvar rascunho}.
\end{remark}

\begin{remark}
\textbf{Imagem recusada em Evidências e fotos}\par
\textbf{Causa provável:} A foto ultrapassa 10 MB, não está em \tech{JPEG}/\tech{PNG}/\tech{WebP} ou já existem 8 fotos anexadas.\par
\textbf{Como resolver:} Reduza o tamanho da imagem ou remova uma das 8 fotos antes de enviar outra.
\end{remark}

\section{Modelos Favoritos e Aplicação de Análise em Lote}

\subsection*{Objetivo}

Permitir que o analista salve uma análise recorrente como \textbf{Modelo Favorito} e aplique esse mesmo preenchimento de uma só vez a várias ocorrências semelhantes da Base de Scrap.\par

\subsection*{Antes de começar}

\begin{itemize}
\item Ter ao menos uma ocorrência já revisada (para salvar um modelo novo) ou um modelo já salvo na lista de \textbf{Modelos Salvos}.
\end{itemize}

\passo[0.50000/0.52121/0.20455/0.43939,0.83182/0.75758/0.88636/0.86364]{completo/21-salvar-modelo-favorito}{Salvar uma análise como modelo favorito}{
Abra uma ocorrência Revisada e clique no ícone de coração para abrir Salvar como modelo favorito.\par
\begin{elementos}
\item \textbf{Nome de referência.} Digite um nome que permita identificar o modelo na lista de Modelos Salvos.
\item \textbf{Salvar.} Grave o modelo para reutilizar o preenchimento em ocorrências semelhantes.
\end{elementos}
}

\begin{remark}[info]
A escolha do modelo, a seleção das ocorrências e a confirmação da aplicação em lote já estão descritas no guia rápido. Consulte o passo~\ref{passo:7} (\pagina{passo:7}).
\end{remark}
\subsection*{Validações e manutenção dos modelos}
\begin{itemize}
\item O painel \textbf{Modelos Salvos} também permite \textbf{Editar modelo} ou \textbf{Excluir} um modelo existente.
\item Antes da confirmação em lote, confira a quantidade de itens selecionados, o \textbf{Tipo de scrap} e a \textbf{Descrição e causa raiz} que serão aplicados.
\item A mensagem de criação dos relatórios confirma o término. Use \textbf{Fechar} para encerrar a janela.
\end{itemize}


\Needspace{160pt}

\subsection*{Campos e regras}

\begingroup
\fontsize{9}{12}\selectfont
\setlength{\LTpre}{4pt}
\setlength{\LTpost}{8pt}
\begin{longtable}{@{}L{3.389cm}L{1.916cm}L{2.652cm}L{6.778cm}@{}}
\toprule
\guiacabtabela \thead{Campo} & \thead{Obrigatório} & \thead{Formato} & \thead{Descrição e regra de preenchimento} \\
\midrule
\endfirsthead
\toprule
\guiacabtabela \thead{Campo} & \thead{Obrigatório} & \thead{Formato} & \thead{Descrição e regra de preenchimento} \\
\midrule
\endhead
\bottomrule
\endlastfoot
\rowcolor{white} \textbf{Nome de referência para este modelo} & Sim & Texto curto & Nome que identifica o modelo dentro da lista \textbf{Modelos Salvos}. \\
\rowcolor{guiaareia!70} \textbf{Caixas de marcação da tabela} & Sim (no modo lote) & Seleção múltipla & Marcam quais linhas da Base de Scrap receberão o preenchimento automático do modelo. \\
\end{longtable}
\endgroup

\subsection*{Mensagens e solução de problemas}

\begin{remark}
\textbf{Ocorrências inativas ou que já possuem relatório serão ignoradas.}\par
\textbf{Causa provável:} Entre os itens marcados na tabela, existe algum registro que já possui relatório vinculado.\par
\textbf{Como resolver:} Mensagem preventiva normal: o sistema aplicará o modelo apenas às ocorrências pendentes e preservará intactas as que já têm relatório.
\end{remark}

\begin{remark}[info]
\textbf{\ensuremath{\checkmark} 3 Relatórios criados com sucesso:}\par
\textbf{O que significa:} O modelo foi aplicado com sucesso a todas as ocorrências válidas selecionadas.\par
\textbf{Próximo passo:} Clique em \textbf{Fechar} para voltar à tabela atualizada.
\end{remark}

\fimdecapitulo
